\section{Catalogue of pools}\label{sec:catalogue}

Table~\ref{tab:catalogue} collects the named uncolorable pools of this version, with the
standard of proof for each minimum.  (The three enlarged integer alphabets of
Section~\ref{sec:real} are not repeated.)  ``Certified'' means the procedure of
Section~\ref{sec:methods} terminated; ``$\le$'' means an uncolorable subset of that size is in
hand and nothing smaller was excluded.  A raw pool is contained in its completion, so a minimum
set of the raw pool bounds the completed pool from above (Lemma~\ref{lem:restrict}); it need not
be a minimum set of it.

\begin{table}[t]
\centering
\caption{Uncolorable pools.  ``raw'': all words over the alphabet; ``compl.'': closure under
third rays of orthogonal pairs.  Letters: $\omega=e^{2\pi i/3}$, $\alpha=(1+\sqrt{-7})/2$,
$\varphi=(1+\sqrt5)/2$, $\psi^3=\psi^2+1$ ($\psi\approx1.4656$), $c=\sqrt[3]{2}$.  Bases: number
of triads inside the minimum (or best) set found.}
\label{tab:catalogue}
\small
\begin{tabular}{llrrrll}
\toprule
Alphabet & Pool & Rays & Pairs & Triads & Minimum KS subset & Bases \\
\midrule
$\{0,\pm1,\pm2\}$ & raw & 49 & 138 & 26 & 31, certified & 17 \\
 & compl. & 109 & 258 & 86 & 31, certified & 17 \\
$\{0,\pm1,\pm\sqrt2\}$ & raw & 49 & 120 & 16 & 33, certified & 16 \\
 & compl. & 145 & 336 & 112 & 33, certified & 20 \\
$\{0,\pm1,\pm\sqrt{-2}\}$ & raw & 49 & 120 & 16 & 33, certified & 16 \\
 & compl. & 145 & 336 & 112 & 33, certified & 20 \\
$\{0,\pm1,\pm\omega,\pm\omega^2\}$ & raw & 57 & 174 & 22 & 33, certified & 14 \\
 & compl. & 165 & 390 & 130 & 33, certified & 14 \\
$\{0,\pm1,\pm i,\pm(1+i)\}$ & raw & 127 & 465 & 51 & 33, certified & 16 \\
 & compl. & 583 & 1521 & 507 & $\le33$ & \\
$\{0,\pm1,\pm\alpha,\pm\alpha^{*}\}$ & raw & 145 & 522 & 42 & 43, certified & 23 \\
 & compl. & 685 & 1746 & 582 & \heegnercompletedmin & \\
$\{0,\pm1,\pm\psi,\pm\psi^2\}$ & raw & 109 & 348 & 40 & 47, certified & 26 \\
 & compl. & \superfourcompleted & \superfourmin \\
$\{0,\pm1,\pm\psi\}$ & compl. & \supercompleted & \supermin \\
$\{0,\pm1,\pm\varphi\}$ & compl. & 205 & 498 & 166 & \goldenmin & \goldenbases \\
$\{0,\pm1,\pm c,\pm c^2\}$ & compl. & \cubiccompleted & \cubicmin \\
\bottomrule
\end{tabular}
\end{table}

\paragraph{Pools that are colorable.}  The raw pools of $\{0,\pm1,\pm\varphi\}$ (49 rays, 10
triads), $\{0,\pm1,\pm\psi\}$ (49 rays, 10 triads) and $\{0,\pm1,\pm c,\pm c^2\}$ (109 rays, 28
triads) are colorable; for these three alphabets uncolorability appears only after completion.
The completed pools of $\{0,\pm1,\pm\sqrt3\}$ and $\{0,\pm1,\pm\sqrt5\}$ each have 133 rays, 282
pairs and 94 triads and are colorable.  (Version~9 printed 145 rays for these two.)

\paragraph{Same field, different alphabet.}  $\sqrt5$ and $\varphi$ generate the same field.  The
completed pool of $\{0,\pm1,\pm\sqrt5\}$ is colorable and that of $\{0,\pm1,\pm\varphi\}$ is not.
This is the plainest evidence that the alphabet, and not the field, is the operative choice.

\subsection{The golden-ratio pool}

The raw pool of $\{0,\pm1,\pm\varphi\}$ is one of the rows of Table~\ref{tab:zerosums}
($1+x-|x|^2=0$): it has more orthogonal pairs than a generic real two-symbol pool (138 against 114)
but the same 10 triads, and it is colorable.  Its completion has 205 rays, 498 pairs and 166
triads and is uncolorable.  Version~9 reported a certified minimum of 52 rays for it.  That
is wrong.  A 51-ray uncolorable subset exists (118 pairs, 28 bases; every single-ray deletion is
colorable); coordinates are in \texttt{golden\_51\_rays.txt}.  The status in this version is:
minimum \goldenmin.

Earlier versions explained the uncolorability of the completed pool by saying that completion
``generates the reciprocal $1/\varphi$''.  That explains nothing, because $1/\varphi$ is already
present projectively: $[\varphi:1:0]=[1:1/\varphi:0]$.  What completion adds is rays, not
scalars, and we have no short account of why these 156 additional rays suffice.

\subsection{Two cubic alphabets}

\paragraph{Supergolden.}  Let $\psi$ be the real root of $\psi^3=\psi^2+1$.  The raw pool of the
four-symbol alphabet $\{0,\pm1,\pm\psi,\pm\psi^2\}$ has 109 rays, 348 pairs and 40 triads and is
uncolorable without any completion.  Its certified minimum is 47 rays (110 pairs, 26 bases).  A
47-ray set over this field was first communicated to us by J.~L.~Gonzalez, extracted from the
completed pool of the two-symbol alphabet $\{0,\pm1,\pm\psi\}$; it refuted the two-mechanism
classification of versions 1--8, because no product of two letters of $\{0,\pm1,\pm\psi\}$ has squared
modulus~2 and the field contains no roots of unity other than $\pm1$.  The relation at work is
the minimal polynomial of $\psi$ itself.  Version~9 answered by adding a third ``mechanism''.  We
now draw the simpler conclusion: any relation among the letters that makes enough inner products
vanish may do, and we have not established a useful finite classification.

\paragraph{Cube root of two.}  With $c=\sqrt[3]2$, the raw pool of $\{0,\pm1,\pm c,\pm c^2\}$ is
colorable (109 rays, 300 pairs, 28 triads).  Its completion has 361 rays, 840 pairs and 280 triads and is
uncolorable.  The smallest uncolorable subset we found has 60 rays (144 pairs, 27 bases) and the
minimum is not certified.  Version~9 reported 58 rays for this pool from the unsound procedure;
we have not reproduced a 58-ray set and do not claim one.

For the completed pools of the two supergolden alphabets (157 and 601 rays) the certifier did not
terminate within its budget; 47 is an upper bound for both, from Gonzalez's set and from the raw
four-symbol pool respectively.

\subsection{Structure of the minimum sets}

\begin{table}[t]
\centering
\caption{Selected KS sets: the minimum sets used in Section~\ref{sec:operational} and the
33-ray set of Conway and Kochen, which is critical but not a minimum set of its pool.  Degree sequences are written
$d^{m}$ for $m$ rays orthogonal to $d$ others.  ``Indisp.'': number of bases whose ``exactly
one'' rule cannot be dropped (all orthogonal-pair exclusions kept) without the set becoming
colorable.  Every set is ray-critical.}
\label{tab:profiles}
\small
\begin{tabular}{lrrrll}
\toprule
Set & Rays & Pairs & Bases & Degree sequence & Indisp. \\
\midrule
$\CK$ & 31 & 71 & 17 & $3^{4}4^{14}5^{8}6^{3}8^{2}$ & 17 of 17 \\
Sch\"utte / Conway--Kochen 33 & 33 & 76 & 20 & $3^{8}4^{13}6^{10}8^{2}$ & 20 of 20 \\
Peres 33 & 33 & 72 & 16 & $4^{30}8^{3}$ & 13 of 16 \\
$\sqrt{-2}$, 33 & 33 & 72 & 16 & $4^{30}8^{3}$ & 13 of 16 \\
Eisenstein 33 & 33 & 78 & 14 & $4^{18}5^{12}8^{3}$ & 13 of 14 \\
$\Q(\sqrt{-7})$, 43 & 43 & 107 & 23 & $4^{16}5^{20}6^{4}8^{2}10^{1}$ & 18 of 23 \\
Supergolden 47 & 47 & 110 & 26 & $3^{14}4^{12}5^{8}6^{6}7^{4}8^{1}9^{2}$ & 26 of 26 \\
Supergolden 47, completed pool$^{\dagger}$ & 47 & 108 & 29 & $3^{16}4^{10}5^{11}6^{3}7^{3}8^{2}9^{2}$ & 29 of 29 \\
Golden 51$^{\dagger}$ & 51 & 118 & 28 & $3^{18}4^{6}5^{12}6^{11}8^{4}$ & 26 of 28 \\
Cube root of two, 60$^{\dagger}$ & 60 & 144 & 27 & $3^{9}4^{24}5^{14}6^{6}7^{3}8^{2}10^{1}12^{1}$ & 26 of 27 \\
\midrule
\multicolumn{6}{l}{$^{\dagger}$ Smallest set found; not certified to be a minimum of its pool.}\\
\bottomrule
\end{tabular}
\end{table}

Table~\ref{tab:profiles} gives the basic invariants.  The three 33-ray sets of Conway--Kochen,
Peres and the Eisenstein alphabet have pairwise non-isomorphic orthogonality graphs (VF2, and
visibly from the degree sequences).  The Eisenstein set is the 33-ray set of
Cabello~\cite{Cabello2025simplest}, or one of its three images in the pool
(Computed fact~\ref{fact:eisenstein}); we have matched the invariants he reports (33 rays, 14
bases); the explicit bijection has not been recorded in the repository.

Whether the 43-ray and 47-ray sets, or the 45-ray family of Theorem~\ref{thm:circle}, appear in
the published catalogues we do not know.  Their graphs differ from the three 33-ray sets and
from $\CK$ by ray count.  A proper comparison against the large automated
catalogues of three-dimensional KS sets has not been done, and the statement of earlier versions
that these configurations had ``not been previously reported'' is withdrawn.
